\section{Introduction}\label{sec:introduction}

The critical FK--Ising model describes the connectivity of the clusters in the
Edwards--Sokal representation of the Ising model. Its planar exploration
interface has the conformally invariant scaling limit $\SLE_{16/3}$.
We prove that this interface limit persists when each lattice coupling is
perturbed by a small independent symmetric two-valued random variable, with
the disorder held fixed while the interface is sampled.

\subsection{The model and the result}

Let $\mathcal E$ be the nearest-neighbor edge set of $\Z^2$, and let
$\xi=(\xi_e)_{e\in\mathcal E}$ be independent random signs, each equally
likely to be $-1$ or $1$. For $0<\eps<1$, put
\begin{equation}\label{eq:random-couplings}
 J_e=1+\eps\xi_e.
\end{equation}
We write $\P_{\mathrm{env}}$ for the law of $\xi$. The critical inverse
temperature is the plus-state magnetization threshold
\[
 \beta_c(\eps)=\inf\bigl\{\beta>0:
       \E_{\mathrm{env}}\langle\sigma_0\rangle^+_{\beta,J}>0\bigr\}.
\]
Here the infinite-volume Gibbs expectation is normalized for each environment
before taking the environmental expectation. This threshold also agrees with
the almost-sure magnetization threshold: positivity propagates between
neighboring vertices for strictly positive ferromagnetic couplings, so its
occurrence is a translation-invariant environmental event. Ergodicity,
monotonicity in $\beta$, and a countable intersection over rational $\beta$
give the equivalence.

Fix a bounded Jordan domain $D\subset\C$ and distinct $a,b\in\partial D$.
A square-lattice domain is a finite filled lattice domain: all its interior
nearest-neighbor edges are present. Boundary collars may separate incidences
at a boundary vertex when recording the order of boundary arcs. Consider
deterministic such domains $G_n=(V_n,E_n)$ at physical mesh $\delta_n\downarrow0$,
with polygonal domains $D_n$ and marked boundary points $a_n,b_n$.
Assume there are orientation-preserving boundary homeomorphisms
$h_n:\partial D\to\partial D_n$ such that
\begin{equation}\label{eq:jordan-approximation}
 \sup_{z\in\partial D}|h_n(z)-z|\longrightarrow0,
 \qquad a_n\longrightarrow a,\quad b_n\longrightarrow b,
\end{equation}
where the marked points correspond in boundary order. All coordinates in
\eqref{eq:jordan-approximation} are physical coordinates, including the
factor $\delta_n$.

Given $\xi$, sample a bond configuration $\omega\in\{0,1\}^{E_n}$ with law
\begin{equation}\label{eq:main-fk-law}
 \phi_{n,\xi}(\omega)
 =\frac1{Z_{n,\xi}}\,
  2^{k_{\mathrm w}(\omega)}
  \prod_{e\in E_n}p_e^{\omega_e}(1-p_e)^{1-\omega_e},
 \qquad p_e=1-e^{-2\beta_c(\eps)J_e}.
\end{equation}
The boundary arc from $a_n$ to $b_n$ is wired and its complementary arc is
free. Thus $k_{\mathrm w}(\omega)$ counts connected components after the
vertices of the wired arc have been identified. The medial exploration
separates the primal open cluster attached to the wired arc from the dual
open cluster adjacent to the free arc. Write $\gamma_n$ for this exploration
in physical coordinates; it is $\delta_n$ times the unit-grid exploration.

Curves are oriented from their first endpoint to their second. We use the
uniform curve metric modulo increasing reparametrization, with constant
pauses immaterial. For representatives $\gamma,\widetilde\gamma:[0,1]\to\C$,
the bounded version of this metric is
\[
 d_{\mathrm{curv}}(\gamma,\widetilde\gamma)
 =1\wedge\inf_{\alpha,\widetilde\alpha}
       \sup_{t\in[0,1]}
       |\gamma(\alpha(t))-\widetilde\gamma(\widetilde\alpha(t))|,
\]
where the infimum is over increasing homeomorphisms of $[0,1]$, followed
by identification of curves at distance zero. Let $d_{\mathrm{BL}}$ denote
the bounded-Lipschitz metric on probability laws on this separable curve
space. The conditional law $\mu_{n,\xi}$ of $\gamma_n$ is a random
probability measure, even though the environment is fixed when an interface
is sampled.

\begin{theorem}\label{thm:main}
There exists $\eps_0>0$ such that for every fixed $0<\eps<\eps_0$, every
bounded marked Jordan domain $(D;a,b)$, and every deterministic filled
square-lattice approximation satisfying \eqref{eq:jordan-approximation},
\[
 d_{\mathrm{BL}}\bigl(\mu_{n,\xi},
       \operatorname{Law}(\SLE_{16/3}(D;a,b))\bigr)
   \xrightarrow[n\to\infty]{\P_{\mathrm{env}}}0.
\]
The constant $\eps_0$ is independent of the domain and its approximation.
\end{theorem}

The disorder strength in Theorem~\ref{thm:main} does not decrease with the
mesh. The conclusion concerns the whole oriented curve, including its
terminal portion near $b$.

\subsection{Background and the source of the difficulty}

The random-cluster representation of Fortuin and Kasteleyn~\cite{FK} and
the coupling of Edwards and Sokal~\cite{ES} connect Ising spin correlations
to random connectivity. At the planar critical point, the fermionic
observables of Smirnov~\cite{Smirnov} and Chelkak--Smirnov~\cite{CS}
give access to conformally invariant boundary-value problems.
Chelkak, Duminil-Copin, Hongler, Kemppainen, and Smirnov~\cite{CDHKS}
proved convergence of critical Ising and FK--Ising interfaces to their
respective SLE limits. Crossing estimates stable under boundary conditions,
as developed by Duminil-Copin, Hongler, and Nolin~\cite{DCHN} and by
Chelkak, Duminil-Copin, and Hongler~\cite{CDH}, and the
curve-compactness framework of Kemppainen--Smirnov~\cite{KS} are central
tools in passing from observables to interfaces.

Random perturbations of the interaction test the stability of this critical
geometry. Harris's criterion~\cite{Harris} identifies the two-dimensional
Ising case as marginal, so the stability question is subtler than in a
strictly irrelevant regime. Perturbative work of Dotsenko--Dotsenko,
Shalaev, and Ludwig investigated weak random bonds and the resulting
logarithmic corrections~\cite{DD,Shalaev,Ludwig,LudwigCorrelations}.
In a recent rigorous treatment, Av\'erous--Mahfouf~\cite{AM} obtain
near-critical random-bond FK crossing estimates with disorder decreasing
with the observation scale for $q>1$. Theorem~\ref{thm:main} concerns a
fixed positive disorder strength and identifies the full interface law.
The present result concerns the specific
symmetric two-valued ferromagnetic disorder in \eqref{eq:random-couplings}.
Its renormalization input is the companion manuscript~\cite{Q}, which
identifies the critical point and constructs a local scale transformation
whose bulk perturbations tend to zero. We state precisely the portions of
that result used here in Section~\ref{sec:inputs}.

Vanishing bulk perturbations alone do not identify FK interface laws.
An explored interface creates a random slit boundary, and conditional
crossing probabilities must be compared in the remaining domain. Moreover,
the available quenched bounds concern fixed physical scales; their
environmental convergence has no rate that would justify a union bound
over all microscopic domains or all possible slit histories. The proof
must preserve these quantifiers when adapting the pure-model arguments.

\subsection{The proof and its reusable ingredients}

Our first step is a partition-function comparison on a fixed polygon.
Near its boundary, bulk symmetry no longer supplies the cancellations
used by the scale transformation. We replace those cancellations with
pure Ising transfer estimates. Along a straight side the relevant
projective distance decays quadratically in the ratio of scales. The
quadratic decay makes the contribution of the boundary cells summable;
at the finitely many corners and changes of boundary condition, a
positive decay exponent suffices. Section~\ref{sec:boundary} proves the
transfer estimates, and Section~\ref{sec:polygon} uses them to extend the
scale transformation to fixed polygon boundaries.

We also need this comparison on polygons lying in a flat covering surface
of a punctured planar domain. Such a polygon has ordinary square-grid
charts, although its projection into the plane need not be one-to-one.
Every lifted edge carries the disorder of its projected edge. The proof
uses independence only within one injective chart and therefore allows
different sheets to carry identical disorder. The pure four-point
observable remains valid on these topological disks, as explained in
Section~\ref{sec:pure}.

Before making these tests depend on an exploration,
Section~\ref{sec:tightness} fixes the environment and controls its curves.
From every mesh subsequence we extract a deterministic further subsequence
on which all polygons in one countable stock have their pure limits and
all annular barriers in a countable dense family hold almost surely.
After fixing such an
environment, those barriers give arm bounds and tightness of the
exploration histories.

The second main step converts fixed tests into conditional tests after
exploration. We puncture near finitely many boundary accesses and lift
the remaining domain to a universal cover. Carefully chosen crosscuts
control its conformal coordinate even near the slit boundary. Small
deletions then confine the central portion of every possible crossing to
a compact part of the cover, where a fixed polygon can test it.
Section~\ref{sec:slits} gives the construction
and the shared-edge comparison. This provides a way to transfer local
partition-function information to adaptive domains while retaining the
original lattice and its disorder.

If a conditional test failed, tightness would yield a convergent sequence
of failing histories. The covering construction would produce one fixed
polygon from the already convergent stock, contradicting its known limit.
Thus the geometry may be chosen after the histories while the final test
remains fixed before the mesh limit. No uniform rate over polygons or
histories is required.

To recover the ordered curve from its driver, Section~\ref{sec:tightness}
also proves a uniform approximation of the tip from positive conformal
height. This uses a
stopping-time version of the small-gate argument in~\cite{KS}.
Finally, Section~\ref{sec:identification} extracts bounded martingale
tests from four-boundary-change probabilities. Two tests identify the
Loewner driver as $\sqrt{16/3}$ times Brownian motion. The tip estimate
and a no-return bound at the target give convergence in the curve metric
of Theorem~\ref{thm:main}.

We adapt the finite-graph drawings, closure comparisons, and martingale
identification strategy of the companion manuscript~\cite{C}, with
their hypotheses stated at the point of use. The boundary extension
and the covering construction are proved here; the homogeneous
probabilistic inputs are replaced by the quenched estimates established
in this paper.
